\SourceLeaf{399}%
\noindent [or, equivalently,
\[
\begin{aligned}
s\sigma(Y\zeta-Z\eta)&=\alpha(t'P+\theta'Q)-a(t'Q+\theta'R)\\
s\sigma(Z\xi-X\zeta)&=\beta(t'P+\theta'Q)-b(t'Q+\theta'R)\\
s\sigma(X\eta-Y\xi)&=\gamma(t'P+\theta'Q)-c(t'Q+\theta'R).
\end{aligned}
\]
Substituting these values of $Y\zeta-Z\eta,\ Z\xi-X\zeta,\ X\eta-Y\xi$ in the expression for $s\frac{df}{du}$, we obtain using 15:]
\[
\begin{aligned}
ss\sigma\frac{df}{du}={}&-(t''\theta'-\theta''t')(PR-QQ)\\
&+(Pt'+Q\theta')\{(-\tfrac12P''t'+Q't'+\tfrac12R'\theta')t'\\
&\hspace{5em}+(\tfrac12R't'+\tfrac12R''\theta')\theta'\}\\
&-(Qt'+R\theta')\{(\tfrac12P't'+\tfrac12P''\theta')t'\\
&\hspace{5em}+(\tfrac12P''t'+Q''\theta'-\tfrac12R'\theta')\theta'\}.
\end{aligned}
\]
[Thus we finally obtain the equation for $\frac{df}{du}$]
\begin{equation}\tag*{[16]}
\left\{\begin{aligned}
&(Pt't'+2Qt'\theta'+R\theta'\theta')\BareRoot{(PR-QQ)}.\frac{df}{du}\\
&\quad=(t'\theta''-\theta't'')(PR-QQ)\\
&\qquad+\tfrac12(Pt'+Q\theta')(-P''t't'+2Q't't'+2R't'\theta'+R''\theta'\theta')\\
&\qquad-\tfrac12(Qt'+R\theta')(P't't'+2P''t'\theta'+2Q''\theta'\theta'-R'\theta'\theta'),
\end{aligned}\right.
\end{equation}
[which contains, apart from $t'=\frac{dt}{du}$, $t''=\frac{ddt}{du^2}$, $\theta'=\frac{d\theta}{du}$, $\theta''=\frac{dd\theta}{du^2}$, only the coefficients $P,Q,R$ of $ss$ and their first derivatives with respect to $t$ and $\theta$.]
\Needspace{9\baselineskip}
\begin{center}\bfseries [5.]\end{center}
\noindent\hspace*{1.2em}To simplify the expression $\left[\text{for }\frac{df}{du}\right]$, introduce [the angle $\chi$ that the tangent to the curve makes with the tangent to the line $\theta=\text{const.}$ at the point of the surface under consideration. Then:]
\begin{equation}\tag*{[17]}
\left\{\begin{aligned}
\operatorname{tang}\chi&=\frac{\theta'\BareRoot{(PR-QQ)}}{Pt'+Q\theta'}\\
\cos\chi&=\frac{Pt'+Q\theta'}{\BareRoot{P}\BareRoot{(Pt't'+2Qt'\theta'+R\theta'\theta')}}\\
\sin\chi&=\frac{\theta'\BareRoot{(PR-QQ)}}{\BareRoot{P}\BareRoot{(Pt't'+2Qt'\theta'+R\theta'\theta')}}.
\end{aligned}\right.
\end{equation}
[Furthermore, let $\omega$ be the angle between the lines $t=\text{const.}$ and $\theta=\text{const.}$,
